%% kapa: exam=1 round=22 year=2011 period=2 form=A subject=appraisal_law subject_ko=부동산관계법규 seq=10 number=50
%% source: Q-Net 감정평가사 기출문제 (cst00309) artlSeq=5036790 file=extracted/1008/5036790/2011년 제22회 감정평가사 1차시험 기출문제/2교시.hwp
%% conversion: kordoc 4.24.51 -> tools/gen1.py (auto); review=proofread-gpt-astra; flags=none
\begin{question}{50}
국토의 계획 및 이용에 관한 법령상 개발밀도관리구역에 관한 설명으로 옳지 않은 것은?
\begin{choices}
  \item 서울특별시장이 서울특별시 중구에서 개발밀도관리구역을 지정하려면 서울특별시 도시계획위원회의 심의를 거쳐야 한다.
  \item 개발밀도관리구역안에서 개발행위를 하려는 자는 기반시설의 설치나 그에 필요한 용지의 확보 등에 관한 계획서를 개발행위허가권자에게 제출하여야 한다.
  \item 주거ㆍ상업 또는 공업지역에서의 개발행위로 기반시설(도시계획시설을 포함한다)의 처리ㆍ공급 또는 수용능력이 부족할 것으로 예상되는 지역 중 기반시설의 설치가 곤란한 지역은 개발밀도관리구역으로 지정될 수 있다.
  \item 용적률의 최대한도가 1천100퍼센트인 상업지역이고 개발밀도관리구역으로 지정된 지역에서는 550퍼센트의 범위에서 용적률을 강화하여 적용한다.
  \item 특별시장ㆍ광역시장ㆍ시장 또는 군수는 개발밀도관리구역을 지정하거나 변경한 경우에는 개발밀도관리구역의 명칭 및 범위, 건폐율 또는 용적률의 강화범위를 지방자치단체의 공보에 게재하여 고시하여야 한다.
\end{choices}
\end{question}
\answer{2}
